\documentclass[11pt,a4paper]{article}
\usepackage[T1]{fontenc}
\usepackage[utf8]{inputenc}
\usepackage{amsmath,amssymb,amsthm}
\usepackage[margin=1in]{geometry}
\usepackage[colorlinks=true,linkcolor=blue,urlcolor=blue]{hyperref}
\theoremstyle{plain}
\newtheorem{theorem}{Theorem}
\newtheorem{lemma}[theorem]{Lemma}
\newtheorem{proposition}[theorem]{Proposition}
\newtheorem{corollary}[theorem]{Corollary}
\theoremstyle{definition}
\newtheorem{remark}[theorem]{Remark}

\title{The empirical recurrence for OEIS A298236, proved by transfer matrix}
\author{Adrian Perez Fontelles\\ \small Independent researcher}
\date{31 August 2026}
\begin{document}
\maketitle

\begin{abstract}
OEIS A298236 counts $4$-column arrays over $\{0,\dots,1\}$ in which every entry has,
among the eight king-move neighbours, a number $equal to$ itself lying in $\{0,2,3,4,5\}$, and it carries an empirical
recurrence of order $47$ contributed by R.~H.~Hardin. It is true. The neighbourhood
reaches one row up and one row down, so the condition on a row involves three consecutive
rows and a window of two rows is a sufficient state; the array is then a walk in a digraph
on $S=256$ vertices. Unlike a plain walk count the boundary rows matter here --- the first
row has nothing above it and the last nothing below --- so the count is
$\mathbf{s}^{\!\top}P^{\,n-2}\mathbf{e}$ for genuine start and end vectors rather than
all-ones. The conjectured recurrence is then annihilated exactly, in integer arithmetic.
\end{abstract}

\noindent\small 2020 Mathematics Subject Classification. 05A15, 05B50, 15B36.\normalsize

\section{The sequence and the conjecture}

OEIS A298236 is ``Number of nX4 0..1 arrays with every element equal to 0, 2, 3, 4 or 5 king-move adjacent elements, with upper left element zero.''. It has offset
$1$ and begins
\[
1, 42, 124, 1156, 8117, 66626, 563826, 4788445,\ \dots
\]
The entry carries the following comment:
\begin{quote}\small
Empirical: a(n) = 8*a(n-1) +32*a(n-2) -170*a(n-3) -756*a(n-4) +1343*a(n-5) +8268*a(n-6) -1088*a(n-7) -52583*a(n-8) -36137*a(n-9) +211652*a(n-10) +261539*a(n-11) -549053*a(n-12) -977458*a(n-13) +949003*a(n-14) +2156328*a(n-15) -1458478*a(n-16) -3452942*a(n-17) +2947713*a(n-18) +5123245*a(n-19) -5969104*a(n-20) -8188927*a(n-21) +10106767*a(n-22) +13662104*a(n-23) -12831855*a(n-24) -19526809*a(n-25) +11999903*a(n-26) +20566810*a(n-27) -10210867*a(n-28) -20355027*a(n-29) +6591071*a(n-30) +17534517*a(n-31) -2850386*a(n-32) -10558997*a(n-33) +518625*a(n-34) +5863076*a(n-35) +420205*a(n-36) -2587827*a(n-37) -463804*a(n-38) +711538*a(n-39) +254620*a(n-40) -205739*a(n-41) -51176*a(n-42) +24022*a(n-43) +3692*a(n-44) -6064*a(n-45) -1992*a(n-46) +624*a(n-47) for n>49
\end{quote}
As of the ``Last modified'' line on the live entry (Jan 16 2018 11:33:14, revision
6) the statement is still recorded as empirical, and nothing on the entry
records it as settled either way.

Written out, the claim is
\begin{equation}\label{eq:conj}
a(n)\;=\;8\,a(n-1) + 32\,a(n-2) - 170\,a(n-3) - 756\,a(n-4) + 1343\,a(n-5) + 8268\,a(n-6) - 1088\,a(n-7) - 52583\,a(n-8) - 36137\,a(n-9) + 211652\,a(n-10) + 261539\,a(n-11) - 549053\,a(n-12) - 977458\,a(n-13) + 949003\,a(n-14) + 2156328\,a(n-15) - 1458478\,a(n-16) - 3452942\,a(n-17) + 2947713\,a(n-18) + 5123245\,a(n-19) - 5969104\,a(n-20) - 8188927\,a(n-21) + 10106767\,a(n-22) + 13662104\,a(n-23) - 12831855\,a(n-24) - 19526809\,a(n-25) + 11999903\,a(n-26) + 20566810\,a(n-27) - 10210867\,a(n-28) - 20355027\,a(n-29) + 6591071\,a(n-30) + 17534517\,a(n-31) - 2850386\,a(n-32) - 10558997\,a(n-33) + 518625\,a(n-34) + 5863076\,a(n-35) + 420205\,a(n-36) - 2587827\,a(n-37) - 463804\,a(n-38) + 711538\,a(n-39) + 254620\,a(n-40) - 205739\,a(n-41) - 51176\,a(n-42) + 24022\,a(n-43) + 3692\,a(n-44) - 6064\,a(n-45) - 1992\,a(n-46) + 624\,a(n-47)
\end{equation}
for all $n>48$.

\section{The array as a walk, with boundary}

For a cell of the array let its \emph{score} be the number of the eight king-move neighbours that lie inside
the array and carry a value $equal to$ the cell's own. The entry's condition is that every cell
has score in $\{0,2,3,4,5\}$. The entry also fixes the top-left entry to be $0$; that is a condition on the first row alone, so it restricts the start vector and nothing else.

A cell's neighbourhood reaches at most one row up and one row down, so the scores of the
cells of a row $s$ are determined by $s$ together with the row above and the row below.
Take as state an ordered pair $(r,s)$ of rows, so $S=256$, and set
\begin{itemize}
\item $\mathbf{s}(r,s)=1$ when every cell of $r$ has score in $\{0,2,3,4,5\}$ \emph{with no row
above $r$} and $r_1=0$, and $0$ otherwise;
\item an edge $(r,s)\to(s,t)$ when every cell of $s$ has score in $\{0,2,3,4,5\}$ computed from
$r$, $s$ and $t$;
\item $\mathbf{e}(r,s)=1$ when every cell of $s$ has score in $\{0,2,3,4,5\}$ \emph{with no row
below $s$}, and $0$ otherwise.
\end{itemize}
Let $P$ be the adjacency matrix of that edge relation.

\begin{lemma}\label{lem:walk}
$a(n)=\mathbf{s}^{\!\top}P^{\,n-2}\mathbf{e}$ for every $n\ge2$.
\end{lemma}

\begin{proof}
An array with $n$ rows is a sequence $r^{(1)},\dots,r^{(n)}$ of rows, and it is counted
exactly when every one of its $n$ rows has all scores in $\{0,2,3,4,5\}$. Read the sequence as the
walk $\bigl(r^{(1)},r^{(2)}\bigr)\to\bigl(r^{(2)},r^{(3)}\bigr)\to\cdots\to
\bigl(r^{(n-1)},r^{(n)}\bigr)$, of length $n-2$. The condition on the first row is
checked by $\mathbf{s}$ at the initial state, because a first row has no row above it; the
condition on row $i$ for $2\le i\le n-1$ is checked by the edge entering
$\bigl(r^{(i)},r^{(i+1)}\bigr)$, which sees exactly $r^{(i-1)}$, $r^{(i)}$ and
$r^{(i+1)}$; and the condition on the last row is checked by $\mathbf{e}$ at the final
state. Each row's condition is therefore tested once and with the correct neighbours, and
summing over all walks gives $\mathbf{s}^{\!\top}P^{\,n-2}\mathbf{e}$.
\end{proof}

\section{The criterion}

Let $q(t)=t^{47}-\sum_i c_i t^{47-i}$ be the characteristic polynomial of
\eqref{eq:conj} and $u_j=\mathbf{s}^{\!\top}P^{j}q(P)\mathbf{e}$.

\begin{lemma}\label{lem:crit}
$a(n)-\sum_i c_i a(n-i)=u_{\,n-2-47}$ whenever $n-47\ge2$. So \eqref{eq:conj}
holds for every $n>t$ exactly when $u_j=0$ for all $j>t-2-47$. Moreover $(u_j)$
satisfies the monic linear recurrence given by the characteristic polynomial of $P$, of
degree $S$, so $S$ consecutive zeros force every later one and the last nonzero $u_j$ pins
$t$ down exactly.
\end{lemma}

\begin{proof}
Substituting Lemma~\ref{lem:walk} and factoring $P^{\,n-2-47}$ out of
$P^{n-2}-\sum_i c_i P^{n-i-2}$ gives the identity, and letting the exponent run gives the
equivalence. For the last part, Cayley--Hamilton gives $\chi(P)=0$ for the characteristic
polynomial $\chi$, so $\chi$ applied to the shift annihilates $(u_j)$; $\chi$ is monic of
degree $S$, which propagates $S$ consecutive zeros forward.
\end{proof}

\section{The computation}

\begin{theorem}
The recurrence \eqref{eq:conj} holds for every $n>48$. This is the statement of
Section~1.
\end{theorem}

\begin{proof}
The vector $q(P)\mathbf{e}\in\mathbb{Z}^{S}$ was formed by $47$ matrix--vector
products in exact integer arithmetic and the $u_j$ evaluated against $\mathbf{s}$ in turn.
They vanish from the index corresponding to $n=48+1$ onwards, and the run was continued
past $S$ consecutive zeros, which by Lemma~\ref{lem:crit} settles every larger index.
\end{proof}

\section{Verification}

Three checks, each independent of the algebra above.

First, the digraph, the start vector and the end vector were built directly from the
entry's stated condition, the one-row count computed separately, and the resulting values
compared against the entry's DATA at every one of the $20$ published indices; they
agree exactly. This is the check that ties the model to the sequence, and it is a sharp one:
getting the boundary rows wrong, or treating them as interior rows, changes the early terms
immediately.

Second, the recurrence \eqref{eq:conj} was evaluated on those published terms in exact
integer arithmetic, with no matrices involved, and vanishes at every index the entry claims.

Third, the vanishing test of Lemma~\ref{lem:crit} was carried past $S=256$ consecutive
zeros rather than to a sampled prefix.

\begin{thebibliography}{9}
\bibitem{oeis} The OEIS Foundation, \emph{The On-Line Encyclopedia of Integer
Sequences}, \url{https://oeis.org}, 2026. Sequence A298236.
\bibitem{stanley} R.~P.~Stanley, \emph{Enumerative Combinatorics}, Volume 1, 2nd ed.,
Cambridge University Press, 2011. (Transfer-matrix method, Section 4.7.)
\bibitem{fs} P.~Flajolet and R.~Sedgewick, \emph{Analytic Combinatorics}, Cambridge
University Press, 2009.
\bibitem{horn} R.~A.~Horn and C.~R.~Johnson, \emph{Matrix Analysis}, 2nd ed., Cambridge
University Press, 2013.
\end{thebibliography}

\end{document}
